\documentclass[11pt]{article}

% Common packages
\usepackage{amsmath}
\usepackage{amssymb}
\usepackage{amsthm}
\usepackage{geometry}
\usepackage{enumitem}
\usepackage{graphicx}
\usepackage{titlesec}
\usepackage{tikz}
\usepackage{pgfplots}
\pgfplotsset{compat=1.18}

% Page setup
\geometry{margin=0.5in}
\setlist[itemize]{topsep=2pt,itemsep=2pt,parsep=0pt}

% --- Load hyperref late, then bookmark AFTER hyperref ---
\usepackage[hidelinks]{hyperref}   % load near the end
\usepackage{bookmark}              % must come after hyperref

% Page setup
\geometry{margin=0.5in}
\setlist[itemize]{topsep=2pt,itemsep=2pt,parsep=0pt}

% Short labels
\newcommand{\thm}{\underline{\textbf{Thm.}} }
\newcommand{\defn}{\underline{\textbf{Def.}} }
\newcommand{\prop}{\underline{\textbf{Prop.}} }
\newcommand{\Gen}{\textsf{Gen}}
\newcommand{\Enc}{\textsf{Enc}}
\newcommand{\Dec}{\textsf{Dec}}

% Title info
\title{CSE 4402 Notes - Cryptography}
\author{Michael Lee}

\begin{document}
\maketitle

\tableofcontents

\newpage

\section{Syllabus / How to Class Works}
\begin{itemize}
    \item 6 homeworks. They are assigned whenever he feels like it.
    \item This class's content is different from MATH 4350
\end{itemize}

\section{Introduction}
\subsection{The Communication Problem}
\begin{itemize}
    \item Typically, we have two characters who are trying to communicate a message (\(m\)) with each other, Alice (\(A\)) and Bob (\(B\)).
    \item However, there is Eve (\(E\)) who is trying to eavesdrop on the communication between \(A\) and \(B\).
    \item We convert \(m\) to a crypt (\(c\)) via a key (\(k\)), send it to \(B\), and then \(B\) decrypts it to get \(m\).
    \item We assume that \(E\) can intercept the crypt (\(c\)) and try to build \(c\) s.t. she cannot decrypt it. \\
    \item \defn Kerchoff's Principle: Assume that \(E\) knows how our system works but doesn't know \(k\).
\end{itemize}

\subsection{Private Key Encryption Scheme}
\textbf{Let the following defintions be true:}
\begin{itemize}
    \item \defn \(m\): message space. \(m \in \mathbb{M}\)
    \item \defn \(k\): key space. \(k \in \mathbb{K}\)
    \item \defn \(c\): cipher. \(c \in \mathbb{C}\)
    \item \defn \(\Gen\): Our key generation procedure (i.e. how we pick \(k \in \mathbb{K}\)) (assume uniformly at random)
    \item \defn \(\Enc\): Our encryption procedure. \(\Enc_k(m) \mapsto c\)
    \item \defn \(\Dec\): decryption procedure. \(\Dec_k(c) \mapsto \{m, \perp\}\)
    \begin{itemize}
        \item Note that \(\perp\) means `null'.
    \end{itemize}
    \item In this class, we display the probability of \(a\) given \(b\) as \(P[b : a]\).
    \item \defn A procedure is `valid' if \(\forall m \in \mathbb{M}, \, P[k \leftarrow \Gen; c \leftarrow \Enc_k(m) : \Dec_k(c) = m] = P_{k,c}[\Dec_k(c) = m] = 1\)
\end{itemize}

\noindent \textbf{Ex. of a valid private key encryption scheme:}
\begin{itemize}
    \item Let \(m\) be 5-digit strings. \(\Gen\) pick \(k\) at random.
    \item Hence, \(P[\text{each key being selected}] = 1/10^5 = 10^{-5}\).
    \item \(\Enc_k(m) : \text{Return \(m\) concatenated with \(k\) (i.e. `\(m||k\)')}\)
    \item For example,
    \[
    \Enc_{10357}(BRIAN) = BRIAN10357
    \]
    \item Obviously, this is NOT very secure, so we should use a more secure scheme.
\end{itemize}

\noindent \textbf{Ex. Caesar Shift:}
\begin{itemize}
    \item \(\mathbb{M}\): finite strings of letters
    \item \(\mathbb{K}\): \(\{0, 1, \ldots, 25\}\)
    \item \(\Enc_k(m)\): Shift every letter in string forward by \(\mathbb{K}\).
    \item \(\Dec_k(c)\): Shift every letter in string backward by \(\mathbb{K}\).
    \item For example,
    \[
    \Enc_{k=3}(PIZZA) = SLCCD
    \]
    \item Obviously, this is NOT very secure, so we should use a more secure scheme.
\end{itemize}

\noindent \textbf{Ex. Substitution Cipher:}
\begin{itemize}
    \item \(\mathbb{M}\): finite strings of letters
    \item \(\mathbb{K}\): permutations of \(\{A, B, C, \ldots, Z\}\)
    \item \(\Enc_k(m)\): \{A, B, C, \ldots, Z\} \(\rightarrow\) \{A, B, C, \ldots, Z\}. Change every instance of letter \(x\) to \(k(x)\).
    \item \(\Dec_k(c)\): change every instance of letter \(y\) to \(k^{-1}(y)\).
    \item Breakable by frequency analysis (count the number of times each letter appears in the ciphertext and compare it to the frequency of each letter in the English language).
    \item Also vulnerable because certain letters are adjacent to other letters (ex. `Q' is often associated with `U').
\end{itemize}

\subsection{\defn Shannon Secrecy}
\begin{itemize}
    \item \defn Shannon Secrecy: Given \(\mathbb{M}\), \(\mathbb{K}\), \(\Gen\), \(\Enc\), \(\Dec\), for any prior distribution, \(D\) on \(\mathbb{M}\),
    \[
        P_{k,m}[m = m_0 | \Enc_k(m) = c] = P_m[m=m_0]
    \]
\end{itemize}

\end{document} 
